\documentclass[11pt]{article}

\usepackage[margin=1in]{geometry}
\usepackage{amsmath,amssymb,amsthm}
\usepackage{enumitem}
\usepackage{hyperref}
\usepackage{xcolor}
\usepackage{booktabs}

\hypersetup{colorlinks=true, linkcolor=blue, urlcolor=blue, citecolor=blue}

\title{TOAE209/TOAH209 -- Design and Analysis of Algorithms\\[4pt]
\large Week 1: Theoretical Questions}
\author{Foreign Trade University}
\date{}

\begin{document}
\maketitle

\section*{Instructions}

Part A contains 22 questions requiring written explanations or short proofs. Part B is a
10-question quiz (true/false and multiple choice). Attempt every question on your own
before consulting Part C (Solutions) at the end of this document.

\section*{Part A: Theory Questions}

\begin{enumerate}[label=\textbf{A\arabic*.}, leftmargin=2.2em]

\item Explain the difference between \emph{worst-case}, \emph{average-case}, and
\emph{best-case} running time. Why does this course (and most of the algorithms
literature) default to worst-case analysis?

\item Prove that Big-O is transitive: if $f(n) = O(g(n))$ and $g(n) = O(h(n))$, then
$f(n) = O(h(n))$.

\item Order the following functions from slowest- to fastest-growing, and justify the
relative order of any two consecutive functions in your list:
\[
  n^2, \quad 2^n, \quad n\log n, \quad \log n, \quad n^3, \quad \sqrt{n}, \quad n, \quad 1.
\]

\item Asymptotic analysis ignores constant factors and lower-order terms (e.g. $3n^2 +
100n + 5 = \Theta(n^2)$). Give a concrete scenario where this abstraction could be
\emph{misleading} in practice -- i.e. where an algorithm with a worse asymptotic bound
could still be faster for the input sizes that actually occur.

\item Compare arrays and singly linked lists in terms of (a) random access, (b) insertion
at the front, (c) insertion at the end, and (d) memory overhead per element. For each of
the four, state which structure is better and why.

\item Explain why a \emph{stack} (LIFO) is the natural data structure for checking
whether brackets in an expression are balanced. What property of nested structures makes
LIFO order the right choice (as opposed to FIFO)?

\item A binary min-heap is stored as an array, where the children of index $i$ are at
indices $2i+1$ and $2i+2$. Explain why both \texttt{push} (insert + sift up) and
\texttt{pop} (remove min + sift down) run in $O(\log n)$ time, in terms of the height of
the tree.

\item Using your answer to A7, explain why \texttt{heapsort} (push all $n$ elements, then
pop them all) runs in $O(n \log n)$ time overall.

\item State the Stable Matching Problem precisely: define a perfect matching, define a
\emph{rogue couple} (blocking pair), and define what it means for a matching to be
\emph{stable}.

\item Prove that the (men-proposing) Gale--Shapley algorithm terminates after at most
$n^2$ iterations of its main loop, where $n = |M| = |W|$.

\item Prove that the matching produced by Gale--Shapley is a \emph{perfect} matching
(every man and every woman ends up matched). Your proof should explain why a woman, once
engaged, can never become \emph{permanently} free again.

\item Prove that the matching produced by Gale--Shapley admits \emph{no rogue couples}
(i.e. it is stable). Your proof should consider an arbitrary pair $(m,w)$ not in the
matching and show it cannot be a rogue couple.

\item Define \emph{valid partner}. State the Men-Optimality Theorem, and explain in your
own words (an informal argument is fine) why the proposing side ends up with their best
possible valid partner.

\item State the Women-Pessimality Theorem. Explain the apparent ``unfairness'' it
describes, and why this is an unavoidable consequence of the Men-Optimality Theorem (not
an independent fact that could go either way).

\item Give a $2\times 2$ instance of the Stable Matching Problem (preference lists for 2
men and 2 women) that has \emph{two} distinct stable matchings. For each of the two
matchings, identify which side (men or women) prefers it.

\item In the Hospital--Residents problem, each hospital $h$ has a capacity
$\text{capacity}(h) > 1$. Explain how the resident-proposing algorithm generalizes
Gale--Shapley, and why the resulting assignment is still ``stable'' (define stability for
this many-to-one setting).

\item With \emph{incomplete} preference lists, some participants may end up unmatched.
Explain how the definition of ``rogue couple'' must be adjusted to account for the
possibility that a participant prefers being unmatched to being matched with someone not
on their list.

\item Karatsuba's algorithm multiplies two $n$-digit numbers using 3 recursive
multiplications of $n/2$-digit numbers, plus $O(n)$ additional work for additions/shifts.
Write down the recurrence $T(n) = \ldots$ and solve it (you may use the Master Theorem,
or unroll the recursion tree) to show $T(n) = \Theta(n^{\log_2 3})$. Approximately what is
$\log_2 3$?

\item Draw (or describe) the recursion tree for \texttt{fib\_naive(5)}. Explain why the
number of calls grows exponentially with $n$, and explain precisely what memoization
changes about this tree.

\item Explain the relationship between the number of recursive calls made by
\texttt{binary\_search\_recursive} on an array of size $n$ and $\log_2 n$. Why does
doubling $n$ only add (roughly) one more call?

\item Insertion sort makes exactly $n-1$ comparisons on an already-sorted array but
$\frac{n(n-1)}{2}$ on a reverse-sorted array. Explain both bounds, and explain what it
means for an algorithm to be ``adaptive'' (i.e. faster on nearly-sorted input).

\item \emph{Discussion (open-ended).} In the Gale--Shapley algorithm, the side that
proposes ends up better off (Men-Optimality) and the side that receives proposals ends up
worse off (Women-Pessimality), \emph{assuming everyone reports their true preferences}.
Suppose a woman could submit a \emph{false} preference list. Could she ever obtain a
better partner than the one she gets under Gale--Shapley with true preferences (assuming
all men still report truthfully)? Discuss informally -- you do not need a formal proof,
but justify your answer with an example or an argument about how the algorithm uses
preference lists.

\end{enumerate}

\newpage
\section*{Part B: Quiz}

\begin{enumerate}[label=\textbf{B\arabic*.}, leftmargin=2.2em]

\item \textbf{(True/False)} If $f(n) = O(n)$, then $f(n) = \Omega(n)$.

\item \textbf{(Short answer)} What is the tightest $\Theta(\cdot)$ bound for
$f(n) = 5n^2 + 100n + 7$?

\item \textbf{(Multiple choice)} A binary min-heap with $n$ elements stored in an array
has height:
\begin{enumerate}[label=(\alph*)]
    \item $O(1)$
    \item $O(\log n)$
    \item $O(n)$
    \item $O(n \log n)$
\end{enumerate}

\item \textbf{(Short answer)} For $n=4$ men and $n=4$ women, what is the \emph{maximum
possible} number of proposals the men-proposing Gale--Shapley algorithm could make
(worst case, using the $n^2$ bound)?

\item \textbf{(Multiple choice)} Which data structure is the natural choice for
simulating round-robin CPU scheduling?
\begin{enumerate}[label=(\alph*)]
    \item Stack (LIFO)
    \item Queue (FIFO)
    \item Binary min-heap
    \item Sorted array
\end{enumerate}

\item \textbf{(True/False)} In Gale--Shapley, once a woman becomes engaged, the
desirability (to her) of her partner is non-decreasing over the rest of the algorithm's
execution -- i.e. she only ever trades up.

\item \textbf{(Short answer)} Karatsuba's algorithm reduces the number of recursive
multiplications needed per level from 4 (grade-school) to how many?

\item \textbf{(Short answer)} For $\text{fib}(20)$ with $F(0)=0, F(1)=1$, the naive
recursive implementation makes 21891 total calls. Roughly how many calls does the
memoized version make (order of magnitude, and exact small multiple of $n$)?

\item \textbf{(True/False)} Every instance of the Stable Matching Problem has exactly one
stable matching.

\item \textbf{(Short answer)} For a sorted array of length $n=15$, what is the maximum
number of recursive calls \texttt{binary\_search\_recursive} makes (i.e.
$\lceil \log_2(n+1) \rceil$)?

\end{enumerate}

\newpage
\section*{Part C: Solutions}

\subsection*{Part A}

\begin{enumerate}[label=\textbf{A\arabic*.}, leftmargin=2.2em]

\item \textbf{Worst-case} running time is the maximum running time over all inputs of a
given size $n$; \textbf{average-case} is the expected running time under some
distribution over inputs; \textbf{best-case} is the minimum. The course defaults to
worst-case because (a) it gives a guarantee that holds for \emph{every} input, with no
assumptions about the input distribution, and (b) average-case analysis requires
committing to a (often unrealistic) probability distribution over inputs, which is harder
to justify and to get right.

\item Since $f(n) = O(g(n))$, there exist $c_1, n_1$ such that $f(n) \le c_1 g(n)$ for all
$n \ge n_1$. Since $g(n) = O(h(n))$, there exist $c_2, n_2$ such that $g(n) \le c_2 h(n)$
for all $n \ge n_2$. Let $n_0 = \max(n_1, n_2)$ and $c = c_1 c_2$. For all $n \ge n_0$,
$f(n) \le c_1 g(n) \le c_1 (c_2 h(n)) = c \cdot h(n)$. Hence $f(n) = O(h(n))$.

\item From slowest- to fastest-growing:
\[
  1 \;\prec\; \log n \;\prec\; \sqrt{n} \;\prec\; n \;\prec\; n\log n \;\prec\; n^2 \;\prec\;
  n^3 \;\prec\; 2^n.
\]
Justifications for consecutive pairs: $1 \prec \log n$ because $\log n \to \infty$ while
$1$ is constant. $\log n \prec \sqrt n$ because $\log n = o(n^\epsilon)$ for any
$\epsilon > 0$. $\sqrt n \prec n$ trivially ($\sqrt n = n^{1/2}$). $n \prec n\log n$ since
$\log n \to \infty$. $n \log n \prec n^2$ since $\log n = o(n)$. $n^2 \prec n^3$ trivially.
$n^3 \prec 2^n$ since any polynomial is $o(2^n)$ (exponentials eventually dominate all
polynomials).

\item Example: an $O(n^2)$ algorithm with very small constants (say running time $\approx
2n^2$ microseconds) versus an $O(n\log n)$ algorithm with large constants (say $\approx
1000 \, n\log n$ microseconds). For $n = 100$: the ``worse'' algorithm takes
$2\cdot 100^2 = 20{,}000$ microseconds, while the ``better'' one takes
$1000 \cdot 100 \cdot \log_2 100 \approx 1000 \cdot 100 \cdot 6.6 \approx 664{,}000$
microseconds. For input sizes that stay in this range (e.g. small fixed-size problems
solved many times), the asymptotically worse algorithm is faster in practice. This is why
real systems often use, e.g., insertion sort for small sub-arrays inside an otherwise
$O(n\log n)$ sort.

\item
\begin{center}
\begin{tabular}{@{}lll@{}}
\toprule
Operation & Better structure & Why \\
\midrule
(a) Random access      & Array       & $O(1)$ index arithmetic vs.\ $O(n)$ pointer walk \\
(b) Insert at front     & Linked list & $O(1)$ new head vs.\ $O(n)$ shift of all elements \\
(c) Insert at end        & Tie / depends & Array: $O(1)$ amortized (with resizing); linked
list: $O(1)$ only with a maintained tail pointer, else $O(n)$ \\
(d) Memory overhead      & Array       & No per-element pointer storage, better cache
locality \\
\bottomrule
\end{tabular}
\end{center}

\item Brackets nest: the bracket that must be closed \emph{next} is always the
\emph{most recently opened} one that is still open. This ``most recent first'' access
pattern is exactly LIFO order. Pushing each opening bracket and popping on each closing
bracket means the top of the stack is always the innermost unmatched opening bracket --
exactly the one the next closing bracket must match. A queue (FIFO) would instead track
the \emph{oldest} unmatched bracket, which is the wrong one.

\item In a binary heap stored as an array, a complete binary tree with $n$ nodes has
height $\lfloor \log_2 n \rfloor$ (each level roughly doubles the number of nodes). Both
sift-up (after \texttt{push}, bubble a node toward the root) and sift-down (after
\texttt{pop}, bubble a node toward the leaves) move along a single path from one level to
an adjacent level, doing $O(1)$ work per level. Since the path length is bounded by the
tree height $O(\log n)$, both operations are $O(\log n)$.

\item \texttt{heapsort} performs $n$ \texttt{push} operations and $n$ \texttt{pop}
operations, each $O(\log n)$ by A7. Total time is $n \cdot O(\log n) + n \cdot O(\log n) =
O(n \log n)$.

\item A \emph{perfect matching} $S$ is a set of $n$ pairs $(m,w)$, $m \in M$, $w \in W$,
such that every man and every woman appears in exactly one pair (a bijection $M
\leftrightarrow W$). A pair $(m,w) \notin S$ is a \emph{rogue couple} (blocking pair) if
$m$ prefers $w$ to his partner in $S$ \emph{and} $w$ prefers $m$ to her partner in $S$. A
matching is \emph{stable} if it is a perfect matching with no rogue couples.

\item Each iteration of the main loop corresponds to one man proposing to one woman he
has \emph{never proposed to before} (a man never repeats a proposal to the same woman).
There are $n$ men, and each has at most $n$ women on his preference list, so the total
number of (man, woman) pairs that could ever be proposed is at most $n \times n = n^2$.
Since each loop iteration consumes one such pair and no pair is reused, the loop runs at
most $n^2$ times.

\item \emph{Key fact:} once a woman becomes engaged, she never becomes free again -- she
can only \emph{switch} to a more-preferred partner (her old partner becomes free, but she
remains engaged, just to someone new). So the set of engaged women is monotonically
non-decreasing over time.

Now suppose, for contradiction, that the algorithm terminates with some man $m$ still
free. Since the loop only stops when no man is both free \emph{and} has unproposed
women left, $m$ must have proposed to (and been rejected by) every one of the $n$ women,
meaning all $n$ women are currently engaged -- to men other than $m$ (since $m$ is free).
But that means at least $n$ men are engaged (one per engaged woman, since engagements are
one-to-one) -- plus $m$ himself is unengaged, giving at least $n+1$ men, contradicting
$|M| = n$. So no man can be left free at termination; since $|M|=|W|=n$ and the matching
is one-to-one, every woman is matched too. Hence the result is a perfect matching.

\item Take any pair $(m,w) \notin S$, where $S$ is the final matching, with $m$ matched to
$w' \ne w$ and $w$ matched to $m' \ne m$. We show $(m,w)$ is not a rogue couple by
considering two cases for whether $m$ ever proposed to $w$:
\begin{itemize}
    \item \textbf{$m$ never proposed to $w$.} Men propose in strictly decreasing order of
    preference, and $m$'s \emph{last} proposal (the one that succeeded) was to $w'$. If
    $m$ never proposed to $w$, it must be because $m$ prefers $w'$ to $w$ (he would have
    reached $w$ eventually otherwise, or stopped at $w'$ which he prefers more). So $m$
    does \emph{not} prefer $w$ to $w'$ -- the first condition for a rogue couple fails.
    \item \textbf{$m$ proposed to $w$ at some point, and was rejected (immediately or
    later displaced).} A woman rejects/displaces a partner only in favor of someone she
    prefers \emph{more}, and once engaged her partner quality is non-decreasing (as
    established in A11). So $w$'s final partner $m'$ is someone $w$ likes at least as
    much as $m$. Hence $w$ does \emph{not} prefer $m$ to $m'$ -- the second condition for
    a rogue couple fails.
\end{itemize}
In both cases, $(m,w)$ is not a rogue couple. Since $(m,w)$ was arbitrary, $S$ has no
rogue couples, i.e. $S$ is stable.

\item A woman $w$ is a \emph{valid partner} of man $m$ if there is \emph{some} stable
matching (not necessarily the one Gale--Shapley produces) in which $m$ and $w$ are paired.
The \textbf{Men-Optimality Theorem} states that in the matching $S^*$ produced by the
men-proposing Gale--Shapley algorithm, every man $m$ is matched to the \emph{best} (most
preferred) woman among all his valid partners.

Informal argument: suppose some man $m$ ended up with a partner $w$ that is \emph{not}
his best valid partner -- i.e. there is a better valid partner $w^*$ that $m$ prefers, and
some stable matching $T$ where $m$ and $w^*$ are paired. Tracing through the algorithm,
$m$ proposes in decreasing order of preference, so he would have proposed to $w^*$ before
$w$. For $m$ to end up with $w$ instead of $w^*$, $m$ must have been rejected by $w^*$ at
some point, which only happens if $w^*$ received a proposal from a man she prefers even
more than $m$. One can show (by a more careful induction over the course of the
algorithm, omitted here) that this would force a contradiction with $T$ being stable --
intuitively, $w^*$'s ``better'' partner in Gale--Shapley would themselves form a rogue
couple with someone in $T$. The full induction is in Kleinberg \& Tardos, Theorem 1.7.

\item The \textbf{Women-Pessimality Theorem} states that in the same matching $S^*$,
every woman is matched with her \emph{worst} valid partner (the least-preferred partner
she has in any stable matching). This seems unfair to the women, but it is a direct
\emph{consequence} of Men-Optimality, not an independent fact: in any \emph{fixed}
matching, if man $m$ is matched to woman $w$, that uses up one ``slot''. If every man
simultaneously gets his \emph{best} valid partner (Men-Optimality), then for any woman
$w$, her partner in $S^*$ must be the \emph{worst} man among those who consider $w$ a
valid partner -- otherwise some man would have to be matched to a partner worse than his
best valid partner, contradicting Men-Optimality for that man. The two theorems are two
views of the same underlying structure (the lattice of stable matchings).

\item Example: men $\{M_1, M_2\}$, women $\{W_1, W_2\}$, where
\[
  M_1: W_1 > W_2, \quad M_2: W_2 > W_1, \quad W_1: M_2 > M_1, \quad W_2: M_1 > M_2
  \quad \text{(a ``rotation'')}.
\]
Both $S_A = \{(M_1,W_1), (M_2,W_2)\}$ and $S_B = \{(M_1,W_2), (M_2,W_1)\}$ are stable (in
each, every man is paired with his \emph{first}-choice woman in $S_A$, or every woman is
paired with her first-choice man in $S_B$ -- check there are no rogue couples in either).
$S_A$ is the men-optimal (and women-pessimal) matching; $S_B$ is the women-optimal (and
men-pessimal) matching. So the \emph{men} prefer $S_A$ and the \emph{women} prefer $S_B$.

\item Each resident proposes down their preference list, one hospital at a time, exactly
as in the one-to-one case. The difference is that a hospital $h$ does not simply hold
``the best resident seen so far'' -- it holds up to \texttt{capacity}$(h)$ residents at
once. When a new resident proposes, if $h$ has fewer than \texttt{capacity}$(h)$ residents
currently held, it accepts; otherwise it compares the new resident to its
\emph{currently-held resident it likes least}, and keeps whichever it prefers, rejecting
the other. \textbf{Stability} here means: no resident $r$ and hospital $h$ (with $r$ not
assigned to $h$) such that $r$ prefers $h$ to their current assignment \emph{and} $h$
either has a free slot or prefers $r$ to one of its currently-assigned residents. The
proof of termination/stability is essentially the same as the one-to-one case, applied
``per slot''.

\item With incomplete lists, a pair $(m,w)$, both currently unmatched (or each preferring
to be unmatched than matched to their current partner, if any), where \emph{both} $m$ and
$w$ find each other acceptable (each appears on the other's list), would form a rogue
couple if $m$ and $w$ both prefer being matched to each other over their current
situation (including the possibility of being unmatched). Formally, ``$m$ prefers $w$ to
his current partner'' must be extended to also cover the case where $m$'s ``current
partner'' is \emph{being unmatched} -- and by the rules of the algorithm, being matched to
anyone on your list is always preferred to being unmatched, so this only matters for
checking that no acceptable pair of currently-unmatched (or mutually-improvable)
participants exists.

\item The recurrence is
\[
  T(n) = 3\,T(n/2) + O(n).
\]
By the Master Theorem (case where $a=3$, $b=2$, and $f(n)=O(n)=O(n^{\log_2 3 - \epsilon})$
is asymptotically smaller than $n^{\log_b a} = n^{\log_2 3}$), the recursive term
dominates and $T(n) = \Theta(n^{\log_2 3})$. Numerically, $\log_2 3 \approx 1.585$, so
$T(n) = \Theta(n^{1.585})$, which is asymptotically better than the grade-school
$\Theta(n^2)$.

(Alternative derivation by unrolling: at depth $k$, there are $3^k$ sub-problems each of
size $n/2^k$, contributing $O(3^k \cdot n/2^k)$ work at that level. Summing over
$k=0,\ldots,\log_2 n$ gives a geometric series dominated by its last term,
$3^{\log_2 n} = n^{\log_2 3}$.)

\item \texttt{fib\_naive(5)} calls \texttt{fib\_naive(4)} and \texttt{fib\_naive(3)};
\texttt{fib\_naive(4)} calls \texttt{fib\_naive(3)} and \texttt{fib\_naive(2)}; and so on.
Crucially, \texttt{fib\_naive(3)} is computed \emph{twice} (once as a child of
\texttt{fib\_naive(5)}'s call to \texttt{fib\_naive(4)}, and once directly), and this
duplication compounds at every level -- the number of calls to compute $F(n)$ satisfies
roughly the same recurrence as $F(n)$ itself, i.e. $C(n) = C(n-1) + C(n-2) + O(1)$, which
grows like $\phi^n$ (the golden ratio to the $n$), i.e. exponentially.

Memoization caches the result of $F(k)$ the first time it's computed. Every subsequent
call to \texttt{fib\_memo}$(k)$ for the \emph{same} $k$ becomes an $O(1)$ cache lookup
instead of a recursive re-computation. Since there are only $n+1$ distinct sub-problems
($F(0), \ldots, F(n)$), and each is computed exactly once (doing $O(1)$ work beyond its
two recursive sub-calls, which are now cache hits), the total work collapses from
exponential to $O(n)$.

\item Each recursive call to \texttt{binary\_search\_recursive} discards \emph{half} of
the remaining search range. Starting from $n$ elements, after $k$ calls the range has
size roughly $n / 2^k$. The recursion stops (range becomes empty) when $n/2^k < 1$, i.e.
$k > \log_2 n$. So the number of calls is $\Theta(\log_2 n)$, specifically
$\lceil \log_2(n+1) \rceil$. Doubling $n$ adds exactly one more halving step needed to
reach an empty range, hence ``roughly one more call'' -- this is the defining
characteristic of logarithmic growth.

\item On an \emph{already-sorted} array, when inserting the $i$-th element (0-indexed),
it is compared only once with its immediate predecessor and found to be already in the
correct position -- so each of the $n-1$ insertions (for $i=1,\ldots,n-1$) costs exactly 1
comparison, for a total of $n-1$. On a \emph{reverse-sorted} array, the $i$-th element
(0-indexed) must be compared with \emph{all} $i$ previously-placed elements before
finding its position (at the very front), so the total is
$1+2+\cdots+(n-1) = \frac{n(n-1)}{2}$.

An algorithm is \textbf{adaptive} if its running time improves on ``nearly sorted'' input
-- insertion sort is adaptive because the number of comparisons is proportional to the
number of \emph{inversions} (pairs of elements out of order) in the input, which is small
for nearly-sorted arrays and large ($\Theta(n^2)$) for reverse-sorted arrays. Many
classic $O(n\log n)$ algorithms (e.g. standard mergesort) are \emph{not} adaptive in this
sense -- they take $\Theta(n\log n)$ regardless of the input order.

\item This is a discussion question, but here is the expected line of reasoning. In the
men-proposing algorithm, women never \emph{act} -- they only passively accept or reject
based on their stated preferences. A woman's outcome is entirely determined by (i) the
true preferences of the men (which determine who proposes to her and in what order), and
(ii) her own stated preference list (which determines which of those proposals she
accepts). It is a known result (related to the Gale--Shapley paper and later work by Dubins
and Freedman, and Roth) that the men-proposing Gale--Shapley mechanism is
\textbf{strategy-proof for the proposing side} (men can never benefit from lying) but
\textbf{not strategy-proof for the receiving side} (a woman can sometimes obtain a better
outcome by mis-stating her preferences -- e.g. by rejecting a proposal from a man she
actually prefers to her eventual Gale--Shapley partner, hoping a better one arrives
later). However, this requires the woman to have information about other participants'
preferences, and is a delicate, risky strategy if her beliefs are wrong (she might end up
unmatched or with someone worse). The key takeaway: \emph{truthfulness is not symmetric}
-- which side benefits from the mechanism's incentive structure depends on which side
proposes.

\end{enumerate}

\subsection*{Part B}

\begin{enumerate}[label=\textbf{B\arabic*.}, leftmargin=2.2em]

\item \textbf{False.} $O(n)$ is an \emph{upper} bound (grows no faster than $n$);
$\Omega(n)$ is a \emph{lower} bound (grows at least as fast as $n$). $f(n) = O(n)$ alone
does not guarantee $f$ grows \emph{at least} linearly -- e.g. $f(n) = 1$ satisfies
$f(n) = O(n)$ but $f(n) \ne \Omega(n)$.

\item $\Theta(n^2)$. The lower-order terms ($100n$, $7$) and the constant factor ($5$)
do not affect the asymptotic growth class.

\item \textbf{(b) $O(\log n)$.} A complete binary tree with $n$ nodes has height
$\lfloor \log_2 n \rfloor$.

\item $n^2 = 4^2 = 16$ proposals (the bound is $n^2$, achieved when $n=4$ gives at most
16).

\item \textbf{(b) Queue (FIFO).} Round-robin processes tasks in arrival order, cycling
each task to the back of the line when its time slice expires -- exactly FIFO behaviour.

\item \textbf{True.} Once engaged, a woman only changes partners when a man she prefers
\emph{more} than her current partner proposes -- so her partner's rank (in her own
preference list) only improves (or stays the same) over time.

\item \textbf{3.} Karatsuba uses the identity $x_1y_0+x_0y_1 = (x_1+x_0)(y_1+y_0) -
x_1y_1 - x_0y_0$ to compute the cross term from the two ``corner'' products plus one extra
product, for 3 total recursive multiplications instead of 4.

\item $O(n)$, specifically $39$ calls for $n=20$ -- roughly $2n$ (each $F(k)$ for
$k=0,\ldots,20$ is computed once, plus a small constant number of bookkeeping calls).

\item \textbf{False.} As shown in A15, an instance can have multiple stable matchings
(e.g. both the men-optimal and women-optimal matchings, when they differ).

\item $\lceil \log_2(16) \rceil = 4$. (Since $n=15$, $n+1=16=2^4$, so
$\lceil \log_2 16 \rceil = 4$.)

\end{enumerate}

\end{document}
